\section{Proof of Theorem~\ref{theorem:semiparametric-master}}
\label{proof:semiparametric-master}

We divide the theorem into two main parts, consistent with the typical
division of asymptotic normality results into a consistency result and a
distributional result. Recall the notation $L(\theta,
\vec{\sigma}) = \E[\loss_{\vec{\sigma},\theta}(Y \mid X)]$
and
$\LtwoP{\vec{\sigma} - \vec{g}}^2
= \E[\ltwo{\vec{\sigma}(Z) - \vec{g}(Z)}^2]$,
where $Z$ has the distribution
that Assumption~\ref{assumption:x-decomposition} specifies.

\subsection{Proof of consistency}

We demonstrate the consistency $\spe \cp u\opt$ in three parts, which we
present as Lemmas~\ref{lemma:super-generic-solutions},
\ref{lemma:vec-sigma-theta-continuity}, and
\ref{lemma:vec-sigma-probability}. The first presents an analogue of
Lemma~\ref{lemma:generic-solutions} generalized to the case in
which there are $m$ distinct link functions, allowing us to characterize
the link-dependent minimizers
\begin{equation*}
  \theta\opt_{\vec{\sigma}} \defeq \argmin_\theta L(\theta, \vec{\sigma})
\end{equation*}
via a one-dimensional scalar on the line $\{t u\opt \mid t \in \R_+\}$.  The
second, Lemma~\ref{lemma:vec-sigma-theta-continuity}, then shows that
$\theta\opt_{\vec{\sigma}} \to u\opt$ as $\vec{\sigma}$ approaches
$\vec{\sigma}\opt$, where the scaling that
$\ltwos{\theta\opt_{\vec{\sigma}}} \to 1$ follows by the normalization
Assumption~\ref{assumption:sp-general}.
Finally, the third lemma
demonstrates the probabilistic convergence $\spe -
\theta\opt_{\vec{\sigma}_n} \cp 0$ whenever $\vec{\sigma}_n \cp
\vec{\sigma}\opt$ in $L^2(P)$.

We begin with the promised analogue of
Lemma~\ref{lemma:generic-solutions}.
\begin{lemma}
  \label{lemma:super-generic-solutions}
  Define the calibration gap function
  \begin{equation}
    \label{eqn:general-h-m-func}
    h_{\vec{\sigma}}(t)
    \defeq \frac{1}{m} \sum_{j = 1}^m
    \E\left[\left(\sigma_j(tZ) - \sigma_j\opt(Z)\right) Z\right].
  \end{equation}
  Then the loss $L(\theta, \vec{\sigma})$ has unique
  minimizer $\theta\opt_{\vec{\sigma}} = t_{\vec{\sigma}} u\opt$ for the
  unique $t_{\vec{\sigma}} \in (0, \infty)$ solving
  $h_{\vec{\sigma}}(t) = 0$. Additionally,
  taking
  $\hessfunc_j(t, z) \defeq
  \sigma_j'(-tz) \sigma_j\opt(z) + \sigma_j'(tz) \sigma_j\opt(-z)$,
  we have
  \begin{equation*}
    \nabla^2 L(t u\opt, \vec{\sigma})
    = \frac{1}{m} \sum_{j = 1}^m \left(
    \E[\hessfunc_j(t, Z) Z^2] u\opt {u\opt}^\top
    + \E[\hessfunc_j(t, Z)] \projperp \Sigma \projperp \right)
    \succ 0.
  \end{equation*}
\end{lemma}
\begin{proof}
	We perform a derivation similar to that we used
	to derive the gap function~\eqref{eqn:h-m-func}, with a few
	modifications to allow collections of $m$ link functions.
	Note that
	\begin{align*}
		L(\theta, \vec{\sigma})
		& = \frac{1}{m} \sum_{j = 1}^m
		\E[\loss_{\sigma_j, \theta}(1 \mid X) \sigma\opt_j(\<X, u\opt\>)
		+ \loss_{\sigma_j,\theta}(-1 \mid X) \sigma\opt_j(-\<X, u\opt\>)],
	\end{align*}
	so that leveraging the usual ansatz that $\theta = t u\opt$, we have
	\begin{align*}
		\nabla L(\theta, \vec{\sigma})
		& = \E\left[\frac{1}{m}
		\sum_{j = 1}^m \left(\sigma\opt_j(-\<X, u\opt\>) \sigma_j(\<\theta, X\>)
		- \sigma\opt_j(\<X, u\opt\>) \sigma_j(-\<\theta, X\>)\right) X \right] \\
		& = \frac{1}{m} \sum_{j = 1}^m
		\E\left[\left(\sigma_j\opt(-Z) \sigma_j(t Z)
		- \sigma_j\opt(Z) \sigma_j(-tZ) Z\right)\right] u\opt \\
		& = h_{\vec{\sigma}}(t) u\opt,
	\end{align*}
	where $h_{\vec{\sigma}}(t) = \frac{1}{m} \sum_{j = 1}^m
	\E[(\sigma_j\opt(-Z) \sigma_j(tZ) - \sigma_j\opt(Z) \sigma_j(-tZ)) Z]$ is
	the immediate generalization of the gap~\eqref{eqn:h-m-func}. We
	now simplify it to the form~\eqref{eqn:general-h-m-func}. Using the
	symmetries $\sigma\opt(z) = 1 - \sigma\opt(-z)$ and $\sigma(z) = 1 -
	\sigma(-z)$ for any $\sigma, \sigma\opt \in \Flink$, we observe that
	\begin{equation*}
		\sigma(tz) (1 - \sigma\opt(z))
		- \sigma(-tz) \sigma\opt(z)
		= \sigma(tz) - (\sigma(tz) + \sigma(-tz)) \sigma\opt(z)
		= \sigma(tz) - \sigma\opt(z),
	\end{equation*}
	and so
	$h_{\vec{\sigma}}(t) = \frac{1}{m} \sum_{j = 1}^m\E[(\sigma_j(tZ) -
	\sigma_j\opt(Z))Z]$.
	Of course, $h_{\vec{\sigma}}(0) = -\frac{1}{m}
	\sum_{j = 1}^m \E[\sigma_j\opt(Z) Z] < 0$, as $\E[Z] = 0$, while
	$\lim_{t \to \infty} h_{\vec{\sigma}}(t) =
	\frac{1}{m} \sum_{j = 1}^m \E[|Z| - \sigma_j\opt(Z) Z] > 0$.
	Then as $h_{\vec{\sigma}}'(t) =
	\frac{1}{m} \sum_{j = 1}^m \E[\sigma_j'(tZ) Z^2] > 0$, there
	exists a unique $t_{\vec{\sigma}} \in (0, \infty)$ satisfying
	$h_{\vec{\sigma}}(t_{\vec{\sigma}}) = 0$.
	
	We turn to the Hessian derivation, where as in the proof of
	Lemma~\ref{lemma:generic-solutions}, we
	write for $\theta = tu\opt$ that
	\begin{align*}
		\nabla^2 L(\theta, \vec{\sigma})
		& = \frac{1}{m} \sum_{j = 1}^m \E\left[
		(\sigma_j'(-t Z) \sigma\opt_j(Z)
		+ \sigma_j'(tZ) \sigma\opt_j(-Z)) Z^2\right] u\opt{u\opt}^\top \\
		& \qquad ~ +
		\frac{1}{m} \sum_{j = 1}^m \E\left[
		(\sigma_j'(-t Z) \sigma\opt_j(Z)
		+ \sigma_j'(tZ) \sigma\opt_j(-Z))\right] \projperp \Sigma
		\projperp,
	\end{align*}
	and so $\nabla^2 L(\theta, \vec{\sigma}) \succ 0$ and $L(\theta,
	\vec{\sigma})$ has unique minimizer $\theta_{\vec{\sigma}}\opt =
	t_{\vec{\sigma}} u\opt$.
\end{proof}

With Lemma~\ref{lemma:super-generic-solutions} serving as the analogue
of Lemma~\ref{lemma:generic-solutions}, we can now show
the continuity of the optimizing parameter $\theta\opt_{\vec{\sigma}}$
in $\vec{\sigma}$:
\begin{lemma}
	\label{lemma:vec-sigma-theta-continuity}
	As $\LtwoP{\vec{\sigma} - \vec{\sigma}\opt} \to 0$,
	we have $\theta\opt_{\vec{\sigma}} - u\opt \to 0$.
\end{lemma}
\begin{proof}
	Via Lemma~\ref{lemma:super-generic-solutions}, it is evidently
	sufficient to show that the solution $t_{\vec{\sigma}}$ to
	$h_{\vec{\sigma}}(t) = 0$ converges to $1$.
	To show this,
	note the expansion
	\begin{equation}
		\label{eqn:expand-h-m}
		m h_{\vec{\sigma}}(t) =
		\sum_{j = 1}^m \left(\E[(\sigma_j(tZ) - \sigma_j\opt(tZ)) Z]
		+ \E[(\sigma_j\opt(tZ) - \sigma_j\opt(Z))Z]
		\right).
	\end{equation}
	We use the following claim, which shows that
	the first term tends to 0 uniformly in $t$ near $1$:
	\begin{claim}
		\label{claim:taco}
		For any $\sigma, \sigma\opt \in \Flink$, we have
		$\sup_{t \in [\half,2]} \LtwoP{\sigma(tZ) - \sigma\opt(tZ)}
		\to 0$ whenever $\LtwoP{\sigma\opt - \sigma} \to 0$.
	\end{claim}
	\begin{proof}
		We use that the density $p(z)$ of $|Z|$ is continuous
		and nonzero on $(0, \infty)$. Take any $0 < M_0 < M_1 < \infty$.
		Then
		\begin{align*}
			\lefteqn{\sup_{t \in [\half, 2]}
				\LtwoP{\sigma\opt(tZ) - \sigma(tZ)}
				=  \sup_{t \in [\half, 2]} \sqrt{\int_0^\infty \prn{\sigma\opt(tz) - \sigma(tz)}^2 p(z) dz}} \\
			& \stackrel{(i)}{\leq}
			\Prb(|Z| \leq M_0) + \Prb(|Z| \geq M_1) +
			\sup_{t \in [\half, 2]} \sqrt{\int_{M_0}^{M_1}
				\prn{\sigma\opt(tz) - \sigma(tz)}^2 p(tz) \cdot
				\frac{p(z)}{p(tz)} dz} \nonumber \\
			& \leq \Prb(|Z| \leq M_0) + \Prb(|Z| \geq M_1) + \sup_{
				\frac{M_0}{2} \le z, z'
				\le 2 M_1}
			\sqrt{\frac{p(z)}{p(z')}}
			\sup_{t \in [\half, 2]} \sqrt{\int_{M_0}^{M_1}
				\prn{\sigma\opt(tz) - \sigma(tz)}^2 p(tz) dz}  ,
		\end{align*}
		where in $(i)$ we use that the link functions $\sigma\opt$ and $\sigma$
		are bounded within $[0,1]$. For any fixed $0 < M_0 \le M_1 < \infty$,
		the ratio $\frac{p(z)}{p(z')}$ is bounded
		for $z, z' \in [\half M_0, 2 M_1]$, and using
		the substitution $tz \mapsto z$ we have
		\begin{align*}
			\sup_{t \in [\half, 2]} \sqrt{
				\int_{M_0}^{M_1}
				\prn{\sigma\opt(tz) - \sigma(tz)}^2 p(tz) dz}
			\leq \sqrt{2} \cdot \LtwoP{\sigma\opt - \sigma}.
		\end{align*}
		We thus have that
		$\LtwoP{\sigma(tZ) - \sigma\opt(tZ)}
		\le K \LtwoP{\sigma - \sigma\opt}
		+ \P(|Z| \not \in [M_0, M_1])$,
		where $K$ depends only on $M_0, M_1$ and the
		distribution of $Z$.
		Take $M_0 \downarrow 0$ and $M_1 \uparrow \infty$.
	\end{proof}
	
	Leveraging the expansion~\eqref{eqn:expand-h-m} preceding
	Claim~\ref{claim:taco} and the claim itself, we see that
	\begin{equation*}
		h_{\vec{\sigma}}(t) = \E[Z^2] \cdot o(1) +
		\frac{1}{m} \sum_{j = 1}^m
		\E[(\sigma_j\opt(tZ) - \sigma_j\opt(Z))Z]
	\end{equation*}
	uniformly in $t \in [\half, 2]$ as $\LtwoP{\vec{\sigma}\opt -
		\vec{\sigma}} \to 0$. The monotonicity of each $\sigma_j\opt$ guarantees
	that if $f_j(t) = \E[(\sigma_j\opt(tZ) - \sigma_j\opt(Z))Z]$, then
	$f'_j(t) = \E[{\sigma_j\opt}'(tZ) Z^2] > 0$, and so $t = 1$ uniquely
	solves $f_j(t) = 0$ and we must have $t_{\vec{\sigma}} \to 1$ as
	$\LtwoP{\vec{\sigma} - \vec{\sigma}\opt} \to 0$.
\end{proof}

Finally, we proceed to the third part of the consistency argument:
the convergence in probability.
\begin{lemma}
	\label{lemma:vec-sigma-probability}
	If $\LtwoP{\vec{\sigma}_n - \vec{\sigma}} \cp 0$,
	then $\spe - \theta_{\vec{\sigma}_n}\opt \cp 0$.
\end{lemma}
\begin{proof}
	By Lemma~\ref{lemma:super-generic-solutions}
	and the assumed continuity of the population Hessian,
	there exists $\delta > 0$ such that
	\begin{equation}
		L(\theta, \vec{\sigma}) \ge L(\theta\opt_{\vec{\sigma}}, \vec{\sigma})
		+ \frac{\lambda}{2} \ltwo{\theta - \theta\opt_{\vec{\sigma}}}^2
		\label{eqn:hessian-lb}
	\end{equation}
	whenever both $\LtwoP{\vec{\sigma} - \vec{\sigma}\opt} \le \delta$ and
	$\ltwo{\theta\opt_{\vec{\sigma}} - \theta} \le \delta$. Applying
	the uniform convergence
	Lemma~\ref{lemma:rademacher-symmetrization}, we see that
	for any $r < \infty$, we have
	\begin{equation*}
		\sup_{\ltwo{\theta} \le r, \vec{\sigma} \in \Flink^m}
		|P_n \loss_{\vec{\sigma},\theta} - L(\theta, \vec{\sigma})| \cp 0.
	\end{equation*}
	For $\delta > 0$, define the events
	\begin{equation*}
		\event_n(\delta) \defeq \left \{
		\ltwo{\theta\opt_{\vec{\sigma}} - u\opt} \le \delta,
		\LtwoP{\vec{\sigma}_n - \vec{\sigma}\opt} \le \delta \right\},
	\end{equation*}
	where
	Lemma~\ref{lemma:vec-sigma-theta-continuity}
	and the assumption that $\LtwoP{\vec{\sigma}_n - \vec{\sigma}\opt}
	\cp 0$ imply that $\P(\event_n(\delta)) \to 1$ for all $\delta > 0$.
	By the growth condition~\eqref{eqn:hessian-lb}
	and uniform convergence
	$P_n \loss_{\vec{\sigma},\theta} - L(\theta,\vec{\sigma}) \cp 0$
	over $\ltwo{\theta} \le r$,
	we therefore have that with probability tending to 1, 
	\begin{equation*}
		\inf_{\ltwos{\theta - \theta\opt_{\vec{\sigma}_n}} = \delta}
		\left\{P_n \loss_{\theta,\vec{\sigma}_n} - P_n \loss_{\theta\opt_{\vec{\sigma}_n},
			\vec{\sigma}_n} \right\}
		\ge \frac{\lambda}{4} \delta^2.
	\end{equation*}
	The convexity of the losses $\loss_{\theta,\vec{\sigma}}$ in $\theta$ and that
	$\spe$ minimizes $P_n \loss_{\theta,\vec{\sigma}_n}$ then
	guarantee the desired convergence $\spe - \theta\opt_{\vec{\sigma}_n} \cp 0$.
\end{proof}

\subsection{Asymptotic normality via Donsker classes}

\providecommand{\G}{\mathbb{G}}

While we do not have $\LtwoP{\vec{\sigma}_n - \vec{\sigma}\opt} =
o_P(n^{-1/2})$, which allows the cleanest and simplest asymptotic normality
results with nuisance parameters
(e.g.~\cite[Thm.~25.54]{VanDerVaart98}), we still expect
$\sqrt{n}(\spe - \theta\opt_{\vec{\sigma}_n})$ to be asymptotically
normal, and therefore, as $\theta\opt_{\vec{\sigma}} = t u\opt$ for some
$t > 0$, the normalized estimators $\spe / \ltwos{\spe}$
should be asymptotically normal to $u\opt$.
To develop the asymptotic normality results, we perform an analysis
of the empirical process centered at the \emph{estimators}
$\spe$, rather than the ``true'' parameter $u\opt$ as would
be typical.

We begin with an expansion. We let $\theta\opt_n =
\theta\opt_{\vec{\sigma}_n}$ for shorthand.
Then as $\P_{n,m} \nabla_\theta \ell_{\spe,
	\vec{\sigma}_n} = 0$ and $\P \nabla_\theta \ell_{\theta\opt_n,
	\vec{\sigma}_n} = 0$, we can derive
\begin{align*}
	\G_{n,m} \nabla_\theta\ell_{\spe, \vec{\sigma}_n} & = \sqrt{n} \prn{\P_{n,m} \nabla_\theta \ell_{\spe, \vec{\sigma}_n} - \P \nabla_\theta \ell_{\spe, \vec{\sigma}_n}} \nonumber \\
	& = \sqrt{n} \prn{\P \nabla_\theta\ell_{\theta\opt_n, \vec{\sigma}_n} - \P \nabla_\theta\ell_{\spe, \vec{\sigma}_n}} \nonumber \\
	& = \sqrt{n} \prn{\nabla_\theta \poploss(\theta\opt_n, \vec{\sigma}_n) - \nabla_\theta \poploss (\spe, \vec{\sigma}_n) } \nonumber \\
	& = \prn{\int_0^1 \nabla_{\theta}^2 \poploss((1-t) \spe + t \theta\opt_n , \vec{\sigma}_n ) d t} \cdot \sqrt{n}(\theta\opt_n - \spe).
\end{align*}
The assumed continuity of $\nabla_\theta^2 \poploss(\theta, \vec{\sigma})$
at $(u\opt, \vec{\sigma}\opt)$ (recall
Assumption~\ref{assumption:sp-general}) then implies that
\begin{align}
	\sqrt{n}(\theta\opt_n - \spe)
	& = \prn{\int_0^1 \nabla_{\theta}^2 \poploss((1-t) \theta\opt_n + t \spe, \vec{\sigma}_n ) d t}^{-1}
	\cdot\, \G_{n,m} \nabla_\theta\loss_{\spe, \vec{\sigma}_n} \nonumber \\
	& =
	\prn{\nabla_\theta^2 \poploss(u\opt, \vec{\sigma}\opt)  + o_P(1)}^{-1}
	\cdot \G_{n,m} \nabla_\theta \loss_{\spe, \vec{\sigma}_n},
	\label{eqn:sp-mle-mid-1}
\end{align}
where we have used the consistency guarantees $\theta\opt_n \cp u\opt, \spe
\cp u\opt$ by Lemmas~\ref{lemma:vec-sigma-theta-continuity}
and~\ref{lemma:vec-sigma-probability} and that $\LtwoP{\vec{\sigma}_n -
	\vec{\sigma}} \cp 0$ by assumption.

The expansion~\eqref{eqn:sp-mle-mid-1} forms the basis of our asymptotic
normality result; while $\spe$ and $\vec{\sigma}_n$ may be data
dependent, by leveraging uniform central limit theorems and the
theory of Donsker function classes, we can show that
$\G_{n,m} \nabla \loss$ has an appropriate normal limit.
To that end,
define the
function classes
\begin{align*}
	\mc{F}_{\delta}
	\defeq \left\{\nabla_\theta \loss_{\theta, \sigma}
	\mid \ltwo{\theta - u\opt} \leq \delta,
	\sigma \in \Flink \right\},
\end{align*}
where we leave the Lipschitz constant $\lipconst$ in $\Flink$ tacit.
By Assumption~\ref{assumption:sp-general} and that $\spe \cp u\opt$, we know
with probability tending to $1$ we have the membership $
\nabla_\theta\ell_{\spe, \vec{\sigma}_n} \in \mc{F}_\delta$. The key
result is then that $\mathcal{F}_\delta$ is a Donsker class:
\begin{lemma}
  \label{lemma:donsker-sp}
  Assume that $\Ep [\ltwo{X}^4] < \infty$.  Then $\mc{F}_\delta$ is a
  Donsker class, and moreover, if $d_{\Flink}((\spe, \vec{\sigma}_n),
  (u\opt, \vec{\sigma}\opt)) \cp 0$, then
  \begin{equation*}
    \G_{n,m} \nabla_\theta\ell_{\spe, \vec{\sigma}_n}
    - \G_{n,m} \nabla_\theta\ell_{u\opt, \vec{\sigma}\opt} \cp 0.
  \end{equation*}
\end{lemma}

Temporarily deferring the proof of Lemma~\ref{lemma:donsker-sp}, let us see
how it leads to the proof of Theorem~\ref{theorem:semiparametric-master}.
Using Lemma~\ref{lemma:donsker-sp} and Slutsky's lemmas in the
equality~\eqref{eqn:sp-mle-mid-1}, we obtain
\begin{align*}
	\sqrt{n}(\spe - \theta_n\opt)
	& = 
	-\prn{\nabla_\theta^2 \poploss(u\opt, \vec{\sigma}\opt)  + o_P(1)}^{-1}
	\cdot \G_{n,m} \nabla_\theta \loss_{u\opt, \vec{\sigma}\opt}
	+ o_P(1) \\
	& \cd \normal\left(0,
	\nabla^2 L(u\opt, \vec{\sigma}\opt)
	\cov\bigg(\frac{1}{m}
	\sum_{j = 1}^m \nabla \loss_{u\opt, \sigma_j\opt}(Y_j \mid X)\bigg)
	\nabla^2 L(u\opt, \vec{\sigma}\opt)\right).
\end{align*}
Calculations completely similar to those we use in the proof of
Theorem~\ref{theorem:multilabel-asymptotics} then give
Theorem~\ref{theorem:semiparametric-master}: because $\P(Y_j = y \mid X = x)
= \sigma_j\opt(y \<u\opt, x\>)$ by assumption,
\begin{equation*}
	\cov\bigg(\frac{1}{m}
	\sum_{j = 1}^m \nabla \loss_{u\opt, \sigma_j\opt}(Y_j \mid X)\bigg)
	= \frac{1}{m^2} \sum_{j = 1}^m
	\cov(\nabla \loss_{u\opt, \sigma_j\opt}(Y_j \mid X))
\end{equation*}
because
$Y_j \mid X$ are conditionally independent, while
\begin{align*}
	\cov(\nabla \loss_{u\opt,\sigma_j\opt}(Y_j \mid X))
	& = \E[\sigma_j\opt(Z)(1 - \sigma_j\opt(Z)) XX^\top] \\
	& = \E[\sigma_j\opt(Z)(1 - \sigma_j\opt(Z)) Z^2]
	u\opt{u\opt}^\top
	+ \E[\sigma_j\opt(Z)(1 - \sigma_j\opt(Z))]
	\projperp \Sigma \projperp.
\end{align*}
When
$\sigma_j = \sigma_j\opt$, the Hessian function $\hessfunc_j$
in Lemma~\ref{lemma:super-generic-solutions}
simplifies to $\hessfunc_j(1, z) = {\sigma_j\opt}'(z)$ as
$\sigma_j\opt$ is symmetric about 0.
We then apply the delta method
as in the proof of Theorem~\ref{theorem:multilabel-asymptotics}.

Finally, we return to the proof of Lemma~\ref{lemma:donsker-sp}.

\paragraph{Proof of Lemma~\ref{lemma:donsker-sp}.}

\newcommand{\ball}{\mathbb{B}}
\newcommand{\Fflip}{\mc{F}_{\mathsf{flip}}}

To prove $\mathcal{F}_\delta$ is Donsker, we show that each coordinate of
$\nabla_{\theta} \loss_{\theta, \sigma} \in \mathcal{F}_\delta$ is, and as
$\nabla_\theta \loss_{\theta, \sigma} (y \mid x) = -y\sigma(-y\<x,
\theta\>)x$, this amounts to showing the coordinate functions $f_{\theta,
	\sigma}^{(i)}(v) = v_i\sigma(-\<v, \theta\>)$ form a Donsker class when
$v$ has distribution $V = YX$. Let
\begin{align*}
	\mathcal{F}_\delta^{(i)}
	& \defeq \left\{ f_{\theta, \sigma}^{(i)}(\cdot) \mid
	\ltwo{\theta - u\opt} \leq \delta, ~
	\sigma \in \Flink \right\},
\end{align*}
so it is evidently sufficient to prove
that $\mathcal{F}_\delta^{(1)}$ forms a Donsker class.

We use bracketing and entropy numbers~\cite{VanDerVaartWe96} to control the
$\mc{F}_\delta^{(i)}$. Recall that for a function class $\mc{F}$, an
\emph{$\epsilon$-bracket} of $\mc{F}$ in $L^q(\P)$ is a collection of
functions $\{(l_i, u_i)\}$ such that for each $f \in \mc{F}$, there exists
$i$ such that $\ell_i \le f \le u_i$ and $\norm{u_i - l_i}_{L^q(\P)} \le
\epsilon$. The \emph{bracketing number} $N_{[\,]}(\epsilon, \mc{F}, L^q(\P))$
is the cardinality of the smallest such $\epsilon$-bracket, and
the \emph{bracketing entropy} is
\begin{align*}
	J_{[\,]}(\mc{F}, L^q(\P))
	\defeq \int_0^\infty
	\sqrt{\log N_{[\,]}(\epsilon , \mathcal{F}, L^q(\Prb))} d \epsilon.
\end{align*}
To show that $\mc{F}_\delta^{(i)}$ is Donsker, it is
sufficient~\cite[Ch.~2.5.2]{VanDerVaartWe96} to show that
$J_{[\,]}(\mc{F}_\delta^{(i)}, L^2(\P)) < \infty$.
%
Our approach to demonstrate that $\mc{F}_\delta^{(1)}$ is Donsker is thus to
construct an appropriate $\epsilon$-bracket of $\mc{F}_\delta^{(1)}$, which
we do
by first covering $\ell_2$-balls in $\R^d$, then for vectors
$\theta \in \R^d$, constructing a bracketing of the induced function class
$\{f_{\theta,\sigma}^{(1)}\}_{\sigma \in \Flink}$, which we combine
to give the final bracketing of $\mc{F}_\delta^{(1)}$.

We proceed with this two-stage covering and bracketing.
Let $\epsilon, \gamma > 0$ be small numbers whose values we determine
later.
Define the $\ell_2$ ball $\ball_2^d = \{x \in \R^d \mid  \ltwo{x} \le 1\}$,
and for any $0 < \epsilon < \delta$, let
$\mc{N}_\epsilon$ be a minimal $\epsilon$-cover of $\delta \ball_2^d$
in the Euclidean norm
of size $N = N(\epsilon, \delta \ball_2^d, \ltwo{\cdot})$,
$\mc{N}_\epsilon = \{\theta_1, \ldots, \theta_N\}$, so that for
any $\theta$ with $\ltwo{\theta} \le \delta$ there exists
$\theta_i \in \mc{N}_\epsilon$ with $\ltwo{\theta - \theta_i} \le \epsilon$.
Standard bounds~\cite[Lemma~5.7]{Wainwright19}) give
\begin{align}
	\label{eqn:ball-covering-number}
	\log N(\epsilon, \delta \ball_2^d, \ltwo{\cdot})
	\le d \log \prn{1 + \frac{2\delta}{\epsilon}}.
\end{align}
For simplicity of notation and to avoid certain tedious negations, we define
the ``flipped'' monotone function family
\begin{equation*}
	\Fflip \defeq \{g : \R \to \R \mid g(t) = \sigma(-t)\}_{
		\sigma \in \Flink}.
\end{equation*}
Now, for any $\theta \in \R^d$, let
$\mu_\theta$
denote the pushforward measure of $\<V, \theta\> = Y\<X, \theta\>$.
For $\theta \in \R^d$, we then let
$\mathcal{N}_{[\,],\gamma, \theta}$ be a minimal
$\gamma$-bracketing of
$\Fflip$
in the  $L^4(\mu_{\theta})$ norm. That is,
for $N = N_{[\,]}(\gamma, \Fflip, L^2(\mu_\theta))$, we have
$\mathcal{N}_{[\,], \gamma, \theta}
= \{(l_{\theta, i}, u_{\theta, i})\}_{i = 1}^{N}$,
and for each $\sigma \in \Flink$,
there exists $i = i(\sigma)$ such that
\begin{align*}
	l_{\theta, i}(t) \leq
	\sigma(-t) \leq
	u_{\theta, i}(t)
	~~ \mbox{and} ~~ \norm{u_{\theta, i} - l_{\theta, i}}_{L^4(\mu_{\theta})}
	\leq \gamma.
\end{align*}
Because elements of $\Fflip \subset \R \to [0, 1]$ are monotone,
\citet[Thm.~2.7.5]{VanDerVaartWe96} guarantee there
exists a universal constant $K < \infty$ such that
\begin{align}
	\label{eqn:monotone-covering-bound}
	\sup_Q \log N_{[\,]}(\gamma, \Fflip, L^4(Q)) \leq \frac{K}{\gamma},
\end{align}
and in particular,
$\log N_{[\,]}(\gamma, \Fflip, L^4(\mu_\theta)) \le \frac{K}{\gamma}$ for
each $\theta \in \R^d$.

With the covering $\mc{N}_\epsilon$ and induced bracketing collections
$\mc{N}_{[\,],\gamma,\theta}$, we now turn to a construction of the
actual bracketing of the class $\mc{F}_\delta^{(1)}$.
For any $\theta \in \R^d$ and bracket $(l_{\theta,i},
u_{\theta, i}) \in \mathcal{N}_{[\,], \gamma, \theta}$, define
the functionals $\what{l}_{\theta,j}, \what{u}_{\theta,j} : \R^d \to \R$ by
\begin{align*}
	\what{l}_{\theta, j}(v)
	& \defeq \hinge{v_1}
	\max\left\{l_{\theta,j}(\<v, \theta\>) - \mathsf{L} \ltwo{v} \epsilon,
	0 \right\}
	- \hinge{-v_1} \min\left\{u_{\theta, j}(\<v, \theta\>) + \mathsf{L} \ltwo{v} \epsilon, 1
	\right\}  , \\
	\what{u}_{\theta, j}(v)
	& \defeq \hinge{v_1} \min\left\{u_{\theta, j}(\<v, \theta\>)
	+ \lipconst \ltwo{v} \epsilon, 1 \right\}
	- \hinge{-v_1} \max\left\{l_{\theta, j}(\<v, \theta\>) -
	\lipconst \ltwo{v} \epsilon, 0 \right\}.
\end{align*}
The key is that these functions form a bracketing
of $\mc{F}_\delta^{(1)}$:
\begin{lemma}
	\label{lemma:bracketing-construction}
	Define the set
	\begin{equation*}
		\mc{B}_{\epsilon,\gamma}
		\defeq
		\left\{ \left(\what{l}_{\theta_i, j},
		\what{u}_{\theta_i, j}\right)
		\mid
		\theta_i \in \mc{N}_\epsilon,
		~ 1 \le j \le N_{[\,]}(\gamma, \Fflip, L^4(\mu_{\theta_i}))\right\}.
	\end{equation*}
	Then
	$\mc{B}_{\epsilon,\gamma}$ is a
	\begin{equation*}
		2 \lipconst \E[\ltwo{X}^4]^{1/2} \cdot \epsilon
		+ \E[\ltwo{X}^4]^{1/4} \cdot \gamma
	\end{equation*}
	bracketing
	of $\mc{F}_\delta^{(1)}$ with cardinality
	at most
	$\log \card(\mc{B}_{\epsilon,\gamma})
	\le \frac{K}{\gamma} + d \log(1 + \frac{\delta}{\epsilon})$.
\end{lemma}
\begin{proof}
	Let $f_{\theta, \sigma}^{(1)}(v) \in \mathcal{F}_\delta^{(1)}$.
	Take $\theta_i \in \mathcal{N}_{\epsilon}$ satisfying
	$\ltwo{\theta - \theta_i} \leq \epsilon$ and
	$(l_{\theta_i, j}, u_{\theta,j}) \in \mc{N}_{[\,], \gamma, \theta_i}$
	such that $l_{\theta_i, j}(t) \leq  \sigma(-t) \leq u_{\theta_i, j}(t)$
	for all $t$, where $\norm{u_{\theta_i,j} - l_{\theta_i,j}}_{L^4(\mu_{\theta_i})}
	\le \gamma$. We first demonstrate the bracketing guarantee
	\begin{align*}
		\what{l}_{\theta_i, j}(v)
		\leq f_{\theta, \sigma}^{(1)}(v) = v_1 \sigma(-\<v, \theta\>)
		\leq \what{u}_{\theta_i, j}(v)
		~~ \mbox{for~all~} v \in \R^d.
	\end{align*}
	For the upper bound, we have
	\begin{align*}
		\lefteqn{f_{\theta, \sigma}^{(1)}(v) = v_1 \sigma(-\<v, \theta\>)}
		\\
		& \stackrel{\mathrm{(i)}}{\leq}
		\hinge{v_1} \min \left\{\sigma(-\<v, \theta_i\>) + \lipconst |\<v, \theta_i - \theta\>|, 1\right\}
		- \hinge{-v_1} \max \left\{\sigma(-\<v, \theta_i\>) - \lipconst |\<v, \theta_i - \theta\>|, 0\right\} \\
		& \stackrel{\mathrm{(ii)}}{\leq} \hinge{v_1}
		\min \left\{\sigma(-\<v, \theta_i\>) + \lipconst \ltwo{v} \epsilon, 1\right\}
		- \hinge{-v_1} \max \left\{\sigma(-\<v, \theta_i\>) - \lipconst \ltwo{v} \epsilon, 0\right\} \\
		& \stackrel{\mathrm{(iii)}}{\leq}
		\hinge{v_1} \min\left\{u_{\theta_i, j}(\<v, \theta_i\>) + \lipconst \ltwo{v} \epsilon, 1 \right\}
		- \hinge{-v_1} \max\left\{l_{\theta_i, j}(\<v, \theta_i\>) - \lipconst \ltwo{v} \epsilon, 0 \right\} \\
		& = \what{u}_{\theta_i, j}(v),
	\end{align*}
	where step (i) follows from the $\lipconst$-Lipschitz continuity of
	$\sigma$, (ii) from the Cauchy-Schwarz inequality and
	that $\ltwo{\theta - \theta_i} \le \epsilon$,
	while step (iii) follows
	by the construction that
	$l_{\theta_i, j}(t) \leq \sigma(-t) \leq u_{\theta_i,
		j}(t)$ for all $t \in \R$.
	Similarly, we obtain the lower bound
	\begin{align*}
		f_{\theta, \sigma}^{(1)}(v) = v_1 \sigma(-\<v, \theta\>) 
		\geq \what{l}_{\theta_i, j}(v),
	\end{align*}
	again valid for all $v \in \R^d$.
	
	The second part of the proof is to bound the distance between the upper
	and lower elements in the bracketing.  By definition,
	$\what{u}_{\theta_i, j} - \what{l}_{\theta_i, j}$ has the pointwise
	upper bound
	\begin{align*}
		\prn{\what{u}_{\theta_i, j}(v) - \what{l}_{\theta_i, j}(v)}^2 &\leq \prn{|v_1| \prn{u_{\theta_i, j}(\<v, \theta_i\>) - l_{\theta_i, j}(\<v, \theta_i\>) + 2 \lipconst \ltwo{v} \epsilon }}^2   .
	\end{align*}
	Recalling that $V = YX$,
	by the Minkowski and Cauchy-Schwarz inequalities, we thus obtain
	\begin{align*}
		\norm{\what{u}_{\theta_i, j}(V) - \what{l}_{\theta_i, j}(V)}_{L^2(\Prb)} & \leq \norm{|V_1| \prn{u_{\theta_i, j}(\<V, \theta_i\>) - l_{\theta_i, j}(\<V, \theta_i\>)} }_{L^2(\Prb)} + \norm{|V_1| \cdot 2 \lipconst \ltwo{V} \epsilon }_{L^2(\Prb)} \nonumber \\
		& \leq \norm{|V_1|}_{L^4(\Prb)} \cdot \prn{\norm{u_{\theta_i, j}(\<V, \theta_i\>) - l_{\theta_i, j}(\<V, \theta_i\>)}_{L^4(\Prb)} + 2\lipconst \epsilon \cdot \norm{\ltwo{V}}_{L^4(\Prb)}  }  .
	\end{align*}
	Noting the trivial bounds $ \norm{|V_1|}_{L^4(\Prb)} \le
	\norm{X}_{L^4(\Prb)} < \infty$ and the assumed bracketing distance
	\begin{align*}
		\norm{u_{\theta_i, j}(\<V, \theta_i\>) - l_{\theta_i, j}(\<V, \theta_i\>)}_{L^4(\Prb)} 	& = \norm{u_{\theta_i, j} - l_{\theta_i, j}}_{L^4(\mu_{\theta_i})}
		\le \gamma,
	\end{align*}
	we have
	the desired bracketing
	distance $\norms{\what{u}_{\theta_i, j} - \what{l}_{\theta_i, j}}_{L^2(\Prb)}
	\leq 2\lipconst \Ep[\ltwo{X}^4]^{1/2} \epsilon
	+ \Ep[\ltwo{X}^4]^{1/4} \gamma$.
	
	The final cardinality bound is immediate via
	inequalities~\eqref{eqn:ball-covering-number}
	and~\eqref{eqn:monotone-covering-bound}.
\end{proof}

Lemma~\ref{lemma:bracketing-construction} will yield the desired entropy
integral bound. Fix any $t > 0$, and note that if we take $\epsilon
=\epsilon(t) \defeq t / (4 \lipconst \E[\ltwo{X}^4]^{1/2})$ and $\gamma =
\gamma(t) \defeq t / (2 \E[\ltwo{X}^4]^{1/4})$, then the set
$\mc{B}_{\epsilon, \gamma}$ is a $t$-bracketing of $\mc{F}_\delta^{(1)}$ in
$L^2$, and moreover, we have the cardinality bound
\begin{align*}
	\log N_{[\,]}(t, \mathcal{F}_{\delta}^{(1)} , L^2(\Prb))
	&
	\leq d \log \left(1 + \frac{2 \delta}{\epsilon(t)}\right)
	+ \frac{K}{\gamma(t)}
	\leq \frac{8\lipconst d \delta \cdot \Ep [\ltwo{X}^4]^{1/2}
		+ 2K \Ep[\ltwo{X}^4]^{1/4}}{t}.
\end{align*}
Additionally, as covering numbers are necessarily integer,
we have $\log N_{[\,]}(t) = 0$
whenever
$t > (8\lipconst d \delta \cdot \Ep [\ltwo{X}^4]^{1/2}
+ 2K \Ep[\ltwo{X}^4]^{1/4}) / \log 2$.
This gives the entropy integral bound
\begin{align*}
	J_{[\, ]}(\mathcal{F}_{\delta}^{(1)}, L^2(\Prb))
	& = \int_0^\infty \sqrt{\log N_{[\,]}(t ,\mathcal{F}_{\delta}^{(1)} ,
		L^2(\Prb))} d t < \infty,
\end{align*}
and consequently (cf.~\cite[Thm.~19.5]{VanDerVaart98} or
\cite[Ch.~2.5.2]{VanDerVaartWe96}), $\mathcal{F}_{\delta}^{(1)}$ is a
Donsker class.  A completely identical argument shows that
$\mathcal{F}_{\delta}^{(i)}, i=2,3,\dots, d$ are Donsker, and so
$\mathcal{F}_\delta$ is a Donsker class, completing the proof of the
first claim in Lemma~\ref{lemma:donsker-sp}.

To complete the proof of Lemma~\ref{lemma:donsker-sp},
we need to show that
\begin{align*}
	\G_{n,m} (\nabla_\theta\loss_{\spe, \vec{\sigma}_n}
	- \nabla_\theta\loss_{u\opt, \vec{\sigma}\opt})
	=
	\frac{1}{m}
	\sum_{j = 1}^m
	\G_n^{(j)} (\nabla_\theta\loss_{\spe, \sigma_{n,j}}
	- \nabla_\theta \loss_{u\opt, \sigma\opt_j})
	\cp 0,
\end{align*}
where $\G_n^{(j)}$ denotes the empirical process on
$(X_i, Y_{ij})_{i=  1}^n$. Notably, because $m$ is finite, it is sufficient
to show that
\begin{equation*}
	\G_n^{(j)} (\nabla_\theta\loss_{\spe, \sigma_{n,j}}
	- \nabla_\theta \loss_{u\opt, \sigma\opt_j})
	\cp 0,
	~~~
	j = 1, \ldots, m.
\end{equation*}
To that end, we suppress dependence on $j$ for notational
simplicity and simply write $\G_n$ and $\loss_{\spe, \sigma_n}$, where
$\LtwoP{\sigma_n - \sigma\opt} \cp 0$.
For any Donsker class $\mc{F} \subset \mc{X} \to \R^d$ and
$\epsilon > 0$,
we have
\begin{equation*}
	\limsup_{\delta \downarrow 0}
	\limsup_{n \to \infty}
	\P\left(\sup_{\LtwoP{f - g} \le \delta}
	\G_n(f - g) \ge \epsilon\right) = 0,
\end{equation*}
(see~\cite[Thm.~3.7.31]{GineNi21}), and so in turn it is sufficient to prove
that
\begin{equation}
	\label{eqn:ltwo-p-convergence-sp}
	\LtwoPbold{\nabla \loss_{\spe, \sigma_n} - \nabla\loss_{u\opt, \sigma\opt}}
	\cp 0.
\end{equation}

To demonstrate the convergence~\eqref{eqn:ltwo-p-convergence-sp}, let $M$ be
finite, and note that for any fixed $\theta, \sigma \in \Flink$ that for $V
= YX$ we have
\begin{align*}
	\LtwoPbold{\nabla \loss_{\theta, \sigma}
		- \nabla \loss_{u\opt, \sigma\opt}}^2
	& = \Ep \brk{\ltwo{V\sigma(-\<V, \theta\>) -
			V\sigma\opt(-\<V,  u\opt\>)}^2} \\
	& \leq \Ep \brk{\ltwo{V}^2 \indic{\ltwo{V} \geq M}}
	+ M^2 \LtwoPbold{\sigma(-\<V, \theta\>) - \sigma\opt(-\<V, u\opt\>)}^2 \\
	& \leq \Ep \brk{\ltwo{X}^2 \indic{\ltwo{X} \geq M}}
	+ 2M^2 \LtwoPbold{\sigma(-\<V, u\opt\>) - \sigma\opt(-\<V, u\opt\>)}^2 \\
	& \qquad\qquad\qquad\qquad\qquad\qquad ~
	+ 2M^2 \LtwoPbold{\sigma(-\<V, u\opt\>) - \sigma(-\<V, \theta\>)}^2
\end{align*}
by the triangle inequality.
As 
\begin{align*}
	\LtwoPbold{\sigma(-\<V, u\opt\>) - \sigma(-\<V, \theta\>)}
	\leq \lipconst \ltwo{\theta - u\opt} \cdot \LtwoPbold{V}
	= \lipconst \ltwo{\theta - u\opt} \cdot \LtwoPbold{X}
\end{align*}
and $\spe - u\opt \cp 0$ and
\begin{align*}
	\LtwoPbold{\sigma_n(-\<V, u\opt\>) - \sigma\opt(-\<V, u\opt\>)}
	= \LtwoPbold{\sigma_n - \sigma\opt} \cp 0
\end{align*}
by assumption,
it follows that for any $\epsilon > 0$ that
\begin{align*}
	\P\left(\LtwoPbold{\nabla_{\theta} \loss_{\spe, \sigma_n}
		- \nabla_{\theta} \loss_{u\opt, \sigma\opt}}
	\ge \E[\ltwo{X}^2 \indic{\ltwo{X} \geq M}] + \epsilon
	\right) \to 0.
\end{align*}
Taking $M \uparrow \infty$ gives the
convergence~\eqref{eqn:ltwo-p-convergence-sp}, completing the proof.


%\appendix
%
%\section{Technical Lemmas}
%
%\begin{lemma}
%	\label{lemma:invert-sums}
%	Let $A, B$ be symmetric matrices with $A + B$ full rank
%	and assume that $AB = BA = 0$. Then
%	$(A + B)^{-1} = A^\dagger + B^\dagger$.
%\end{lemma}

